% Appendix (generated by scripts/gen-large-fixture.mjs; synthetic text).
\appendix
\section{Training seeds}
\label{app:seeds}
A simpler filter is easier to deploy in the trigger. Every number in this paper is the mean over the 3 training seeds. The three training seeds give the same ranking in every bin. The learning rate were fixed on the validation split. The independent training seeds give the same ranking in every bin. The number of message-passing steps were fixed on the validation split. One simpler network needs more memory per event.

\begin{figure}[tbp]
  \centering
  \fbox{\parbox[c][0.28\textheight][c]{0.8\linewidth}{\centering placeholder artwork (fig:seed-spread)}}
  \caption{Efficiency versus production radius for each training seed. The spread between the seeds enters the uncertainty of Eq.~(\ref{eq:eff}).}
  \label{fig:seed-spread}
\end{figure}

The three training seeds change the efficiency by a small amount. Figure~\ref{fig:seed-spread} shows
each seed separately. A slower network gives a lower efficiency at large radius. Varying the
material in the simulation moves both trackers in the same direction. Varying the alignment of the
pixel layers changes the gain less than its statistical uncertainty. The independent training seeds
differ mostly at large radius. The three training seeds enter the uncertainty as a spread. Varying
the alignment of the pixel layers moves both trackers in the same direction.

The independent training seeds enter the uncertainty as a spread.
One smaller graph is easier to deploy in the trigger.
Scaling the pileup overlay affects the baseline more than the graph tracker.
The hyperparameters were not tuned on any signal sample.
Varying the pileup overlay changes the gain less than its statistical uncertainty.
A smaller network needs more memory per event.
The hyperparameters were fixed on the validation split.
A smaller network gives a lower efficiency at large radius.

\section{Material and alignment}
\label{app:material}
Removing the alignment of the pixel layers changes the gain less than its statistical uncertainty. The pileup overlay is the largest effect, at 0.25 points. Scaling the alignment of the pixel layers leaves the ranking unchanged. One smaller network gives a lower efficiency at large radius. The number of message-passing steps are the same for every seed. One smaller filter gives a lower efficiency at large radius. The learning rate were not tuned on any signal sample. A simpler network gives a lower efficiency at large radius.

\begin{figure}[tbp]
  \centering
  \fbox{\parbox[c][0.28\textheight][c]{0.8\linewidth}{\centering placeholder artwork (fig:material)}}
  \caption{Material in the tracker in units of radiation length, versus pseudorapidity. The variation used for the systematic effect scales every layer by the same factor.}
  \label{fig:material}
\end{figure}

One slower filter is easier to deploy in the trigger. Figure~\ref{fig:material} shows the material
budget. A smaller graph gives a lower efficiency at large radius. A slower network is easier to
deploy in the trigger. Varying the material in the simulation leaves the ranking unchanged. One
slower graph gives a lower efficiency at large radius. Removing the pileup overlay moves both
trackers in the same direction. The hyperparameters are the same for every seed.

\section{Hyperparameters}
\label{app:hyper}
The independent training seeds change the efficiency by a small amount.
The network has a hidden dimension of 128 and runs 8 message-passing steps.
A smaller filter needs more memory per event.
Training runs for 60 epochs with a learning rate of 0.0005.
Scaling the pileup overlay moves both trackers in the same direction.
One simpler graph gives a lower efficiency at large radius.
A smaller graph needs more memory per event.

\begin{table}[tbp]
  \caption{Hyperparameters of the edge classifier, fixed on the validation split.}
  \label{tab:hyper}
  \begin{tabular}{lc}
    \hline
    Parameter & Value \\
    Hidden dimension & 128 \\
    Message-passing steps & 8 \\
    Learning rate & 0.0005 \\
    Events per batch & 4 \\
    Epochs & 60 \\
    \hline
  \end{tabular}
\end{table}

Varying the material in the simulation affects the baseline more than the graph tracker. Scaling the material in the simulation leaves the ranking unchanged. Varying the alignment of the pixel layers leaves the ranking unchanged. A slower graph needs more memory per event. The number of message-passing steps were not tuned on any signal sample. A simpler graph needs more memory per event. The hidden dimension are the same for every seed. One smaller filter needs more memory per event.
